\chapter{Безопасность}

Контейнеры дают изоляцию, но не являются границей безопасности уровня
виртуальной машины. Ошибка в Dockerfile, лишняя привилегия или пароль
в слое образа сводят защиту на нет. В главе~--- практические меры,
которые закрывают большинство рисков без усложнения эксплуатации.

\section{Модель угроз}

\begin{center}
\small
\begin{tabularx}{\textwidth}{@{}XX@{}}
\toprule
\textbf{Риск} & \textbf{Мера} \\
\midrule
уязвимости в базовом образе и зависимостях & минимальный образ,
регулярная пересборка, сканирование \\
лишние привилегии процесса & \cmd{USER}, \cmd{--cap-drop},
\cmd{no-new-privileges} \\
секреты в образе или истории сборки & \file{.dockerignore},
\flag{--secret}, менеджер секретов \\
доступ к сокету Docker & не монтировать в контейнеры, использовать
proxy или rootless \\
заражённый образ из реестра & официальные источники, digest, подписи,
SBOM \\
атака через общее ядро & актуальное ядро, rootless, gVisor/Kata
для недоверенного кода \\
исчерпание ресурсов & лимиты памяти, CPU, числа процессов \\
открытые наружу порты & публикация только нужного, \cmd{127.0.0.1},
внутренние сети \\
\bottomrule
\end{tabularx}
\end{center}

\begin{defbox}
\textbf{Принцип минимальных привилегий}: контейнер получает только
те права, файлы и сетевые доступы, которые нужны приложению для
работы. Всё остальное отключено по умолчанию.
\end{defbox}

\section{Запуск без root}

По умолчанию процесс в контейнере работает от \cmd{root}. Если
приложение скомпрометировано, атакующий получает root внутри
контейнера~--- и, при ошибке в конфигурации, путь к хосту.

\begin{code}[language=dockerfile]
FROM python:3.13-slim

RUN groupadd --gid 10001 app \
 && useradd --uid 10001 --gid app --create-home app

WORKDIR /app
COPY --chown=app:app . .
RUN pip install --no-cache-dir -r requirements.txt

USER 10001:10001
CMD ["python", "app.py"]
\end{code}

\begin{itemize}
  \item \cmd{USER} в Dockerfile задаёт пользователя по умолчанию.
  \item \cmd{--user 10001:10001} в \cmd{docker run} переопределяет его.
  \item В Kubernetes удобнее числовые UID: они не зависят от
        \file{/etc/passwd} образа.
\end{itemize}

\begin{code}[language=bash]
$ docker run -d --name api \
    --user 10001:10001 \
    --read-only \
    --tmpfs /tmp:size=32m \
    --cap-drop=ALL \
    --security-opt=no-new-privileges:true \
    --pids-limit=200 \
    --memory=512m --cpus=1 \
    api:1.0
\end{code}

\begin{itemize}
  \item \flag{--read-only} делает корневую файловую систему только для
        чтения; запись возможна в тома и \cmd{tmpfs}.
  \item \flag{--cap-drop=ALL} убирает все Linux capabilities; при
        необходимости добавьте одну: \flag{--cap-add=NET\_BIND\_SERVICE}.
  \item \flag{--security-opt=no-new-privileges:true} запрещает
        программы с setuid повышать права.
  \item \flag{--pids-limit} защищает от fork-бомб.
\end{itemize}

\begin{warnbox}
\flag{--privileged} отключает большинство ограничений: контейнер
получает доступ к устройствам, все capabilities и зачастую может
влиять на хост. Не используйте его, если точно не понимаете,
зачем он нужен.
\end{warnbox}

\section{User namespaces и rootless}

В обычном режиме \cmd{root} внутри контейнера~--- это тот же UID~0,
что и на хосте, но с ограниченными возможностями. User namespaces
отображают внутренние UID на непривилегированные внешние.

\begin{code}[language=bash]
{
  "userns-remap": "default"
}
\end{code}

Это фрагмент \file{/etc/docker/daemon.json}; после перезапуска демон
переотобразит UID и GID. Побочный эффект: существующие тома и образы
могут стать недоступны, а права на привязанные каталоги придётся
согласовывать.

Rootless-режим (раздел~3.4.2) идёт дальше: весь демон работает
от обычного пользователя. Для рабочих станций и CI это хороший
компромисс между удобством и безопасностью.

\begin{notebox}
Rootless и user namespaces не отменяют уязвимости ядра, но снижают
последствия: процесс с UID~0 внутри контейнера не является root
на хосте.
\end{notebox}

\section{Секреты}

Секреты нельзя хранить в образе: слои сохраняются, \cmd{docker history}
и сканеры их видят. Основные правила:

\begin{itemize}
  \item не передавайте секреты через \cmd{ARG}~--- они попадают
        в историю и кэш;
  \item не копируйте \file{.env} и ключи в образ~--- исключайте
        их в \file{.dockerignore};
  \item передавайте секреты при запуске через переменные окружения,
        файлы, смонтированные только для чтения, или менеджер
        секретов;
  \item не логируйте значения секретов, маскируйте их в приложении.
\end{itemize}

Compose поддерживает секреты-файлы:

\begin{code}[language=composeyaml]
services:
  api:
    image: api:1.0
    secrets:
      - db_password

secrets:
  db_password:
    file: ./secrets/db_password.txt
\end{code}

Внутри контейнера секрет появится как файл
\file{/run/secrets/db\_password} с правами только для чтения.
В Swarm тот же синтаксис использует зашифрованное хранилище менеджера.

\begin{warnbox}
Переменные окружения видны процессам контейнера и могут попасть
в логи, дампы и диагностические отчёты. Для паролей и ключей
предпочтительнее файлы или специализированные хранилища (Vault,
облачные secret manager'ы), а не \cmd{ENV}.
\end{warnbox}

\section{Безопасность образов}

\subsection{Минимальная база}

Чем меньше в образе пакетов и утилит, тем меньше поверхность атаки.
Порядок предпочтений: \cmd{scratch} или distroless~$\rightarrow$
\cmd{alpine}/slim~$\rightarrow$ полный дистрибутив.

\begin{code}[language=dockerfile]
FROM gcr.io/distroless/static-debian12:nonroot
COPY --from=build /out/app /app
ENTRYPOINT ["/app"]
\end{code}

\subsection{Сканирование}

\begin{code}[language=bash]
$ docker scout cves api:1.0
$ trivy image --severity HIGH,CRITICAL api:1.0
$ docker history --no-trunc api:1.0
$ docker buildx build --sbom=true --provenance=true --push -t api:1.0 .
\end{code}

Сканируйте не только финальный образ, но и базовый: уязвимость
в \cmd{python:3.13-slim} перейдёт во все производные образы.

\subsection{Проверка происхождения}

\begin{code}[language=bash]
$ cosign verify --certificate-identity-regexp '.*' \
    --certificate-oidc-issuer https://token.actions.githubusercontent.com \
    ghcr.io/acme/api:1.0
\end{code}

Подпись и проверка подписи позволяют разрешать к запуску только
образы, собранные доверенным конвейером. Политики внедряются
в Kubernetes (Kyverno, OPA Gatekeeper) или в admission-контроллеры
реестра.

\section{Сокет Docker}

Монтирование \file{/var/run/docker.sock} в контейнер даёт ему
возможность управлять демоном, а значит~--- запускать привилегированные
контейнеры и монтировать каталоги хоста. Это фактически root.

\begin{code}[language=bash]
# Опасно: контейнер получает управление Docker на хосте
$ docker run -v /var/run/docker.sock:/var/run/docker.sock my-ci-image
\end{code}

Безопасные альтернативы:

\begin{itemize}
  \item \textbf{proxy}~--- \cmd{docker-socket-proxy} разрешает только
        нужные методы API;
  \item \textbf{отдельный демон} в rootless-режиме или в виртуальной
        машине для CI;
  \item \textbf{SSH-контекст} вместо проброса сокета;
  \item \textbf{BuildKit-демон} с ограниченным набором операций.
\end{itemize}

\section{Сеть}

\begin{itemize}
  \item Публикуйте только необходимые порты и по возможности
        на \cmd{127.0.0.1}.
  \item Не используйте \flag{--network host} без необходимости.
  \item Для внутренних сервисов создавайте сеть с флагом
        \flag{--internal}: у неё нет выхода наружу.
  \item Не смешивайте в одной сети сервисы разных уровней доверия.
\end{itemize}

\begin{code}[language=bash]
$ docker network create --internal backend
\end{code}

\begin{notebox}
\flag{--internal} запрещает маршрутизацию наружу, но не заменяет
firewall на хосте. Контейнеры во внутренней сети всё равно доступны
друг другу, поэтому сегментируйте приложение по сетям.
\end{notebox}

\section{Ограничения ресурсов}

Лимиты~--- не только производительность, но и защита от атак
на истощение.

\begin{code}[language=bash]
$ docker run -d --name api \
    --memory=512m --memory-swap=512m \
    --cpus=1 \
    --pids-limit=200 \
    --ulimit nofile=1024:2048 \
    api:1.0
\end{code}

В Compose те же параметры задаются в \cmd{deploy.resources.limits}
и \cmd{pids\_limit} (в зависимости от версии Compose).

\section{Аудит и соответствие}

\begin{itemize}
  \item \textbf{Docker Bench for Security}~--- скрипт проверки
        конфигурации демона и хоста по рекомендациям CIS.
  \item \textbf{Регулярные обновления}~--- ядро, Docker Engine,
        базовые образы.
  \item \textbf{Логи аудита}~--- кто запускал контейнеры и с какими
        привилегиями.
  \item \textbf{Политики}~--- запрет \flag{--privileged},
        обязательные лимиты, разрешённые реестры.
\end{itemize}

\section{Чек-лист перед production}

\begin{enumerate}
  \item Процесс работает не от root, UID зафиксирован.
  \item Корневая файловая система только для чтения, запись~--- в тома.
  \item Сброшены capabilities, включён \cmd{no-new-privileges}.
  \item Заданы лимиты памяти, CPU и числа процессов.
  \item Секретов нет в образе, истории и переменных окружения.
  \item Образ собран из минимальной базы и просканирован.
  \item Базовые образы зафиксированы по версии или digest.
  \item Наружу опубликованы только нужные порты.
  \item Сокет Docker не смонтирован в приложение.
  \item Настроены healthcheck, логи и обновление базовых образов.
\end{enumerate}

\begin{summarybox}
\begin{itemize}
  \item Контейнер не равен виртуальной машине по уровню изоляции.
  \item Минимальные привилегии: non-root, read-only, cap-drop,
        лимиты.
  \item Секреты~--- только во время запуска, не в слоях образа.
  \item Образы нужно сканировать, подписывать и фиксировать по digest.
  \item Сокет Docker равносилен root; используйте proxy, rootless
        или отдельный демон.
\end{itemize}
\end{summarybox}

\section*{Вопросы для самопроверки}
\begin{selfcheck}
  \item Почему контейнер под root опаснее, чем процесс под обычным
        пользователем?
  \item Что делает \flag{--read-only} и куда тогда писать временные
        файлы?
  \item Почему секрет в \cmd{ARG} считается утёкшим?
  \item Чем опасен смонтированный \file{/var/run/docker.sock}?
  \item Какие меры снижают риск запуска заражённого образа?
  \item Зачем нужен \flag{--pids-limit}?
\end{selfcheck}
